\begin{frame}{첫 스텝에서는 ``정직한 첫 걸음''}
  \begin{block}{핵심 성립}
    $L^{\mathrm{IS}}$ 의 그래디언트는 \textbf{평범한 policy gradient
    $\nabla_\theta\log\pi_\theta\, A_t$ 에 계수 $r_t$ 가 곱해진
    형태}다. 데이터를 방금 모은 시점
    ($\theta=\theta_{\mathrm{old}}$)에서는 $r_t=1$ 이므로
    \textbf{10.3절 Actor-Critic의 그래디언트와 정확히 같다}.
  \end{block}
  \vskip0.5em
  \begin{itemize}
    \item 즉, 확률비로 목적함수를 다시 쓴 순간
      \kb{``옛 데이터를 재사용해도 첫 스텝에서는 정직한 policy
      gradient''}가 자동으로 성립 — 이것이 확률비가 ``재사용의
      열쇠''인 이유.
    \item $\theta$ 가 $\theta_{\mathrm{old}}$ 에서 멀어지면 계수
      $r_t$ 가 1에서 벗어나 이 동등성이 \textbf{약해진다} —
      ``멀어지는 정도''를 막는 것이 11.2절의 클리핑.
  \end{itemize}
\end{frame}
